% Purpose: Academic introduction to earthquake mechanics, wave propagation physics, regional seismotectonics, SMINO network, classical vs. ML picking, and research roadmap.
% Status: Draft; under academic review.
% Source-of-truth: OGS/src/, SMINO documentation, literature citations in lib.bib.

\label{chap:Introduction}

\section{Seismological Foundations and Wave Propagation Physics}
\label{sec:wave_physics}

Earthquakes represent transient mechanical disturbances generated by the sudden release of accumulated elastic strain energy within the Earth's lithosphere. The lithosphere behaves as a brittle-elastic medium under tectonic differential stress; when shear stresses acting along a pre-existing fault plane exceed the frictional resistance governed by Byerlee's law or rate-and-state friction, dynamic shear failure nucleates and ruptures along the fault surface. This catastrophic dislocation radiates transient elastodynamic stress waves that propagate through the heterogeneous interior and along the free surface of the Earth.

In a continuous, isotropic, linearly elastic medium characterized by mass density $\rho(\mathbf{x})$ and Lamé elastic moduli $\lambda(\mathbf{x})$ and $\mu(\mathbf{x})$, the infinitesimal elastodynamic equation of motion expressing conservation of linear momentum is formulated as:
\begin{equation}
  \rho(\mathbf{x}) \frac{\partial^2 u_i}{\partial t^2}(\mathbf{x}, t) = \frac{\partial \sigma_{ij}}{\partial x_j}(\mathbf{x}, t) + f_i(\mathbf{x}, t),
  \label{eq:momentum_conservation}
\end{equation}
where $\mathbf{u}(\mathbf{x}, t) = [u_1, u_2, u_3]^{\mathrm{T}}$ denotes the displacement vector field, $\boldsymbol{\sigma}(\mathbf{x}, t)$ is the symmetric Cauchy stress tensor, and $\mathbf{f}(\mathbf{x}, t)$ represents external body force distributions per unit volume (such as equivalent double-couple source representations of fault rupture). Under the generalized Hooke's law for isotropic elasticity, the constitutive stress-strain relation is given by:
\begin{equation}
  \sigma_{ij} = \lambda \delta_{ij} \epsilon_{kk} + 2\mu \epsilon_{ij},
  \label{eq:hooke_law}
\end{equation}
where $\delta_{ij}$ is the Kronecker delta, $\epsilon_{ij} = \frac{1}{2}\left(\frac{\partial u_i}{\partial x_j} + \frac{\partial u_j}{\partial x_i}\right)$ is the infinitesimal strain tensor, and Einstein summation convention over repeated indices is implied. Substituting the constitutive law~\eqref{eq:hooke_law} into the equation of motion~\eqref{eq:momentum_conservation} for a homogeneous medium yields the Navier-Cauchy elastodynamic wave equation:
\begin{equation}
  \rho \frac{\partial^2 \mathbf{u}}{\partial t^2} = (\lambda + \mu)\nabla(\nabla \cdot \mathbf{u}) + \mu \nabla^2 \mathbf{u} + \mathbf{f}.
  \label{eq:navier_cauchy}
\end{equation}

By invoking Helmholtz's decomposition theorem, any smooth vector displacement field $\mathbf{u}(\mathbf{x}, t)$ can be partitioned uniquely into an irrotational (curl-free) scalar potential $\phi(\mathbf{x}, t)$ and a solenoidal (divergence-free) vector potential $\boldsymbol{\psi}(\mathbf{x}, t)$:
\begin{equation}
  \mathbf{u} = \nabla \phi + \nabla \times \boldsymbol{\psi}, \quad \text{with} \quad \nabla \cdot \boldsymbol{\psi} = 0.
  \label{eq:helmholtz}
\end{equation}
Substituting decomposition~\eqref{eq:helmholtz} into the homogeneous equation of motion~\eqref{eq:navier_cauchy} completely decouples the wavefield into two independent wave equations propagating through the medium at distinct characteristic phase velocities:
\begin{equation}
  \frac{\partial^2 \phi}{\partial t^2} = \alpha^2 \nabla^2 \phi, \quad \frac{\partial^2 \boldsymbol{\psi}}{\partial t^2} = \beta^2 \nabla^2 \boldsymbol{\psi},
  \label{eq:decoupled_wave_equations}
\end{equation}
where $\alpha$ and $\beta$ represent the longitudinal compressional (Primary, P) and transverse shear (Secondary, S) wave propagation velocities, respectively:
\begin{equation}
  \alpha = \sqrt{\frac{\lambda + 2\mu}{\rho}}, \quad \beta = \sqrt{\frac{\mu}{\rho}}.
  \label{eq:wave_velocities}
\end{equation}

The physical characteristics and particle motion trajectories of these seismic phases govern the observational records captured by modern seismological sensor systems:
\begin{enumerate}[label=(\alph*)]
  \item \textbf{Primary (P) Waves}: Compressional or dilatational longitudinal waves that propagate via volumetric deformation ($\nabla \cdot \mathbf{u} \neq 0, \nabla \times \mathbf{u} = \mathbf{0}$). Particle motion is strictly parallel to the direction of wave propagation. Because $\lambda > 0$ and $\mu > 0$, $\alpha > \beta$ holds universally, ensuring that P waves represent the fastest seismic phase and form the first onset arrival observed on seismograms. P waves propagate through solids, fluids, and gases.
  \item \textbf{Secondary (S) Waves}: Equivoluminal transverse shear waves that propagate via shear deformation without volume change ($\nabla \cdot \mathbf{u} = 0, \nabla \times \mathbf{u} \neq 0$). Particle displacement occurs perpendicular to the ray propagation vector. In a Cartesian coordinate system aligned with the incidence plane, S waves are decomposed into vertically polarized shear motion (SV) and horizontally polarized shear motion (SH). Because fluids have zero shear modulus ($\mu = 0$), S waves propagate exclusively through solid elastic media.
\end{enumerate}

In an ideal Poissonian solid ($\lambda = \mu$, equivalent to Poisson's ratio $\nu = 0.25$), the kinematic velocity ratio satisfies:
\begin{equation}
  \frac{\alpha}{\beta} = \sqrt{\frac{\lambda + 2\mu}{\mu}} = \sqrt{3} \approx 1.732.
  \label{eq:poisson_ratio}
\end{equation}
In complex geological crustal formations, lithological variations, microcrack porosity, and fluid saturation cause the observed $V_{\mathrm{P}}/V_{\mathrm{S}}$ ratio to vary typically between $1.65$ and $1.90$.

At discontinuous geological interfaces and the Earth's stress-free surface ($\sigma_{i3}|_{z=0} = 0$), mode conversion and interference between body waves generate \textbf{surface waves}. These propagate horizontally along the crustal waveguide with amplitudes decaying exponentially with depth:
\begin{itemize}
  \item \textbf{Rayleigh Waves}: Resulting from constructive P-SV interference at the free boundary, exhibiting retrograde elliptical particle motion in the vertical-radial propagation plane, with phase velocities $c_{\mathrm{R}} \approx 0.92\beta$.
  \item \textbf{Love Waves}: Resulting from constructive SH wave trapping within low-velocity surface layers overlying a higher-velocity half-space, exhibiting purely horizontal transverse motion perpendicular to the propagation path.
\end{itemize}

Modern digital seismological observatories record these vector ground motions along three mutually orthogonal spatial components: Vertical (Z), North-South (N), and East-West (E). Observational stations deploy high-gain broadband seismometers (e.g., measuring velocity $\dot{\mathbf{u}}(t)$ over flat bandwidths from $0.01$ to $50\,\mathrm{Hz}$) and strong-motion accelerometers (e.g., measuring ground acceleration $\ddot{\mathbf{u}}(t)$ without clipping during high-amplitude near-source shaking). Data are continuously digitized at high sampling rates ($100\,\mathrm{Hz}$) and packaged into standard miniSEED (MSEED) containers, with instrumental transfer functions, calibration poles, zeros, and geodetic station coordinates formally cataloged in StationXML schemas.

\section{Seismotectonic Context of Northeastern Italy and the 1976 Legacy}
\label{sec:tectonics}

Northeastern Italy, embracing the Friuli Venezia Giulia (FVG) and Veneto regions, represents one of the most tectonically active and seismically hazardous sectors of the entire European Alpine arc. The geodynamic evolution of the region is governed by the active collision and northward indentation of the Adriatic microplate (Adria) into the stable Eurasian plate. Geodetic measurements from continuous GNSS networks indicate that the Adria microplate rotates counterclockwise around an Eulerian pole in the Western Po Plain, producing active crustal shortening across the Eastern Southern Alps at rates of $1.5\text{--}3.0\,\mathrm{mm/yr}$ in a NNW-SSE orientation.

Structurally, Northeastern Italy constitutes a complex tectonic junction where two major collisional orogenic belts converge:
\begin{enumerate}
  \item The \textbf{Eastern Southern Alps}: A south-verging thrust-and-fold retro-wedge characterized by E-W to WNW-ESE striking thrust faults and ramp-anticlines (including the Periadriatic Lineament, the Fella-Sava fault, the Maniago thrust, and the Buia-Tricesimo thrust system) that accommodate crustal shortening through north-dipping blind and exposed thrusting.
  \item The \textbf{External Dinarides}: A dextral strike-slip structural system extending from Slovenia into Eastern Friuli, dominated by NW-SE trending transcurrent fault structures such as the Idrija and Ravne fault networks.
\end{enumerate}

This tectonic configuration results in concentrated intraplate seismicity characterized by shallow hypocenters (typically $5\text{--}15\,\mathrm{km}$ depth), predominantly compressive thrust-faulting mechanisms with subordinate strike-slip kinematics, and recurrent destructive earthquake sequences.

The defining modern seismic episode in this region is the catastrophic \textbf{1976 Friuli earthquake sequence}. On May 6, 1976, at 20:00:12 UTC, a major earthquake of moment magnitude $M_{\mathrm{w}}~6.4$ ($M_{\mathrm{L}}~6.5$) nucleated near Gemona del Friuli, producing epicentral intensities of $I_0 = \mathrm{IX}\text{--}\mathrm{X}$ on the Mercalli-Cancani-Sieberg (MCS) scale. The mainshock ruptured blind thrust faults along the Southalpine front. The sequence was compounded in September 1976 by two devastating aftershocks on September 11 ($M_{\mathrm{w}}~5.9$) and September 15 ($M_{\mathrm{w}}~6.0$). In total, the 1976 Friuli seismic crisis claimed 989 lives, injured over 2,800 people, left over 100,000 homeless, and destroyed entire medieval villages across the Tagliamento and Fella river basins.

The tragedy of the 1976 disaster catalyzed a profound transformation in European seismology, earthquake engineering, civil protection organization, and disaster response legislation. Crucially, it led the Regional Government of Friuli Venezia Giulia to mandate the establishment of a dedicated permanent seismological monitoring infrastructure. In 1977, the \acf{OGS} established the \acf{CRS} in Udine and Trieste. Since its inception, OGS CRS has maintained continuous, high-resolution seismic monitoring across Northeastern Italy, serving as the official technical-scientific body supporting the regional Civil Protection agencies of Friuli Venezia Giulia and Veneto.

\section{The OGS-SMINO Monitoring Network Infrastructure}
\label{sec:smino_network}

Over four decades of technological development, the initial sparse analog network deployed after the 1976 earthquake evolved into the modern \acf{SMINO} (\cite{SMINO}). The SMINO monitoring infrastructure represents a state-of-the-art, multi-parametric, real-time observation system covering Friuli Venezia Giulia, Veneto, and surrounding border areas.

The regional monitoring infrastructure integrates:
\begin{itemize}
  \item A dense permanent backbone of digital three-component broadband and very-broadband seismometers (principally Streckeisen STS-2, Nanometrics Trillium, and Guralp sensors) installed in vaults, tunnels, and boreholes to minimize cultural and atmospheric noise.
  \item High-dynamic-range strong-motion accelerometers (Kinemetrics Episensor) collocated at broadband sites and deployed across urban centers and critical engineering infrastructures to ensure on-scale recordings of strong ground motions.
  \item Cross-border data exchange with neighboring Italian, Austrian, and Slovenian institutions, creating a unified transboundary seismic monitoring domain.
\end{itemize}

In operational practice, seismic monitoring in Northeastern Italy is not conducted in isolation; wavefield propagation respects no administrative boundaries. As described in Chapter~\ref{chap:Datasets}, the complete regional monitoring aperture analyzed in this study encompasses \textbf{266 stations} originating from 18 distinct regional, national, and international networks (including OGS, \ac{INGV}, GeoSphere Austria, the Slovenian Environment Agency ARSO, the Swiss Seismological Service SED, and the Italian National Accelerometric Network RAN). The continuous recording from this dense collective array provides the empirical foundation for our machine learning and computational investigations.

\section{Classical Automated Picking and Association Limitations}
\label{sec:classical_picking}

Transforming continuous, multi-channel ground motion velocity streams into discrete, scientifically verified earthquake catalogs requires identifying the exact onset times of P- and S-wave phase arrivals (phase picking) and grouping arrivals originating from the same physical hypocentral rupture (phase association).

For over four decades, automated seismological processing has relied on classical heuristic algorithms:
\begin{enumerate}
  \item \textbf{Short-Time Average over Long-Time Average (\ac{STA/LTA})}: Introduced in foundational form by \textcite{STALTA}, the method calculates the ratio of energy or amplitude within a short sliding time window ($T_{\mathrm{STA}}$) to that within a preceding long time window ($T_{\mathrm{LTA}}$) over a characteristic function $\mathit{CF}(t)$:
  \begin{equation}
    R(t) = \frac{\mathit{STA}(t)}{\mathit{LTA}(t)} = \frac{\frac{1}{T_{\mathrm{STA}}}\int_{t-T_{\mathrm{STA}}}^t \mathit{CF}(\tau)\,\mathrm{d}\tau}{\frac{1}{T_{\mathrm{LTA}}}\int_{t-T_{\mathrm{LTA}}}^t \mathit{CF}(\tau)\,\mathrm{d}\tau},
    \label{eq:stalta}
  \end{equation}
  where $\mathit{CF}(t)$ is commonly defined as $y(t)^2$ or $y(t)^2 + C \dot{y}(t)^2$. A pick trigger is declared when $R(t) \ge \eta_{\mathrm{on}}$, and deactivated when $R(t) \le \eta_{\mathrm{off}}$.
  \item \textbf{Autoregressive Akaike Information Criterion (\ac{AIC})}: Applied to seismic phase timing by \textcite{LEONARD1999247}, the AIC models the seismogram before and after an arbitrary split point $k$ as two distinct stationary autoregressive processes. The minimum of the AIC function identifies the optimal transition between background noise and seismic onset:
  \begin{equation}
    \mathit{AIC}(k) = k \log(\hat{\sigma}_1^2) + (N - k) \log(\hat{\sigma}_2^2),
    \label{eq:aic}
  \end{equation}
  where $\hat{\sigma}_1^2$ and $\hat{\sigma}_2^2$ are the variance estimators of the prediction errors for intervals $[1, k]$ and $[k+1, N]$, respectively.
\end{enumerate}

Despite their enduring historical utility, classical algorithms suffer from fundamental physical and algorithmic limitations:
\begin{itemize}
  \item \textbf{Vulnerability to Noise and False Triggers}: STA/LTA is highly sensitive to non-stationary transient noise, cultural vibrations, electrical glitches, and meteorological disturbances, causing excessive false trigger rates in densely populated or industrial valleys.
  \item \textbf{Emergent and Low-SNR Degradation}: Onsets of small-magnitude earthquakes ($M_{\mathrm{L}} < 2.0$) or emergent arrivals traveling through attenuative, heterogeneous crustal paths exhibit low signal-to-noise ratios (\ac{SNR}). STA/LTA triggers systematically late, while AIC minimization frequently settles into local minima created by noise fluctuations.
  \item \textbf{Phase Misidentification}: Neither STA/LTA nor AIC natively distinguishes between compressional (P) and shear (S) phases without ad-hoc, brittle horizontal-to-vertical amplitude ratios or polarization filters. S-wave arrivals arriving within the high-amplitude, reverberating P-wave coda are regularly missed or misclassified.
  \item \textbf{Combinatorial Association Bottlenecks}: Grid-search associators that cluster picks by evaluating travel-time residuals across spatial velocity grids exhibit exponential computational scaling $\mathcal{O}(N_{\mathrm{picks}}^k)$ with pick volume. When spurious noise picks enter the pipeline, association algorithms suffer combinatoric explosion or form hallucinated, unphysical hypocenters.
\end{itemize}

Because of these limitations, operational seismological observatories have historically required dedicated human analysts to inspect, refine, or manually pick phase arrivals for every candidate event. This manual intervention introduces severe operational latency (often hours to days), creates subjective timing inconsistencies between analysts, and establishes an insurmountable operational ceiling that forces observatories to discard microseismic events below completeness thresholds ($M_{\mathrm{c}} \approx 1.5$).

\section{The Big Data Era and Deep Learning in Seismology}
\label{sec:dl_seismology}

Over the past decade, the rapid expansion of permanent networks, dense nodal geophone arrays, and real-time transboundary telemetry has propelled observational seismology into the era of ``Big Data.'' Continuous seismic archives at regional and national data centers ingest tens of gigabytes per day and routinely accumulate terabytes to petabytes of continuous waveform data annually. Analyzing these continuous, multi-channel data streams exhaustively is beyond human manual processing capacity.

This data deluge has coincided with revolutionary advances in \acf{DL}. Deep neural networks possess the unique ability to learn rich, non-linear, hierarchical feature representations directly from continuous, uncurated sensor time series without requiring manual feature engineering. Applied to seismic waveforms, deep convolutional and attention-based architectures learn the multi-scale spatio-temporal signatures of wave propagation, dispersion, impedance contrasts, and noise textures.

Key architectural breakthroughs include:
\begin{itemize}
  \item \textbf{PhaseNet} (\cite{PhaseNet}): A deep 1D Convolutional Neural Network (\ac{CNN}) based on the U-Net architecture. PhaseNet employs an encoder-decoder topology with skip connections, processing three-component waveforms through successive downsampling and upsampling convolutional stages to produce continuous probability distributions for P-wave arrivals, S-wave arrivals, and noise.
  \item \textbf{Earthquake Transformer (\ac{EQTransformer})} (\cite{EQTransformer}): A sophisticated multi-task deep network that integrates 1D convolutional layers, bidirectional \ac{LSTM} recurrent units, and multiple multi-head self-attention mechanisms. EQTransformer models both local waveform geometries and long-range temporal dependencies across 60-second windows, simultaneously outputting detection probabilities for earthquake presence, P-wave arrival, and S-wave arrival.
\end{itemize}

To overcome software fragmentation and provide standardized interfaces for benchmarking, \textcite{SeisBench} introduced \textbf{SeisBench} (\cite{SeisBench}), an open-source platform offering unified Python APIs for machine learning in seismology. SeisBench provides standardized data loaders, benchmark training datasets, and pre-trained weights for state-of-the-art architectures, dramatically enhancing the reproducibility of deep-learning seismic workflows. In this thesis, SeisBench serves as the core foundation for evaluating pre-trained phase pickers across diverse training corpora.

\section{The March 27, 2024 Socchieve Earthquake Sequence}
\label{sec:socchieve_sequence}

To test the operational efficacy, robustness, and scientific fidelity of deep-learning pipelines under demanding real-world conditions, this research focuses on an intense seismic sequence that struck Friuli in early 2024.

On \textbf{March 27, 2024, at 21:19:37 UTC} (22:19:37 local time), a magnitude $M_{\mathrm{L}}~4.6 \pm 0.3$ ($M_{\mathrm{w}}~4.5$) earthquake occurred in the Carnic Prealps, approximately $5\,\mathrm{km}$ south-southwest of Socchieve, Udine province (hypocentral coordinates: Latitude $46.3583^\circ\,\mathrm{N}$, Longitude $12.8080^\circ\,\mathrm{E}$, focal depth $10.22\,\mathrm{km}$). The earthquake was strongly felt across the entire Friuli Venezia Giulia and Veneto regions, as well as in neighboring Austria and Slovenia, causing localized structural damage, masonry cracking, and rockfalls in the epicentral area.

This earthquake constitutes the \textbf{most energetic tectonic event to occur in Friuli since the 1976--1980 crisis}. It initiated a vigorous aftershock sequence comprising hundreds of minor earthquakes and microseismic events that decayed over subsequent weeks and months. The spatial proximity of dense seismic instrumentation, combined with the presence of complex Alpine topography, variable focal mechanisms, and realistic operational cultural noise, renders the Socchieve sequence a pristine natural testbed for stress-testing machine learning detection, phase association, and hypocentral relocation against an expert analyst ground truth.

\section{Research Objectives and Scope of the Dissertation}
\label{sec:objectives}

The overarching scientific objective of this doctoral dissertation is to construct, optimize, deploy, and rigorously validate a complete, automated, procedural machine learning pipeline for earthquake pick detection, phase association, and event characterization, specifically engineered for scalable execution on \ac{HPC} infrastructures.

The key research questions addressed in this thesis are:
\begin{enumerate}
  \item \textbf{Domain Adaptation and Generalization}: How effectively do deep-learning phase pickers (PhaseNet and EQTransformer) trained on large external datasets (INSTANCE in Italy, STEAD globally, SCEDC in California) generalize to the complex geological and noise environment of Northeastern Italy? Does regional training data (INSTANCE) outperform massive but geographically foreign corpora (SCEDC, STEAD)?
  \item \textbf{Associator Scaling and Performance}: How do advanced phase association algorithms—specifically the Bayesian Gaussian Mixture Model Associator (\ac{GaMMA}) and the high-throughput octree-based associator (\ac{PyOcto})—compare in association recall, computational efficiency, and robustness against false-positive pick contamination?
  \item \textbf{Physically Bounded Performance Validation}: How can model-generated earthquake catalogs be compared against manual analyst reference catalogs without introducing spatiotemporal bias or duplicate assignment artifacts?
  \item \textbf{HPC Scalability and Operational Viability}: Can the entire processing workflow—from continuous multi-client FDSN waveform ingestion to deep inference, association, and relocation—be parallelized across modern GPU-accelerated supercomputing architectures (such as the CINECA Leonardo Tier-0 system) to support near-real-time operational monitoring and rapid historical archive reprocessing?
\end{enumerate}

\section{Dissertation Structure and Chapter Roadmap}
\label{sec:roadmap}

The remainder of this dissertation is organized into eight interconnected chapters, followed by a dedicated Data and Resources statement and technical appendices:
\begin{itemize}
  \item \textbf{Chapter~\ref{chap:Datasets} (Datasets)}: Details the operational case study data acquired across Northeastern Italy (March 20 to June 20, 2024, 266 stations, 380 relocated OGS reference events with 7,129 picks) and reviews the four benchmark training corpora (INSTANCE, STEAD, SCEDC, and Original sets).
  \item \textbf{Chapter~\ref{chap:data_engineering} (Data Engineering)}: Presents the software architecture of the pipeline, detailing the central \texttt{ogsconstants} configuration, weight-to-uncertainty mappings, automated FDSN multi-client waveform downloading, multi-format legacy catalog parsing (HPL, DAT, TXT, PUN), catalog aggregation, geodetic math, and visualization frameworks.
  \item \textbf{Chapter~4 (Computational Pipeline)}: Formulates the complete procedural workflow, detailing continuous waveform preprocessing, deep-learning phase picking inference, Bayesian (\ac{GaMMA}) and octree (\ac{PyOcto}) phase association, 1D \ac{NonLinLoc} probabilistic relocation, and local magnitude ($M_{\mathrm{L}}$) inversion.
  \item \textbf{Chapter~5 (Pipeline Evaluation Framework)}: Develops the formal evaluation methodology, defining bipartite graph maximum-weight matching (\ac{BPGMA}) for both picks and events, the construction of domain-specific confusion matrices, and the mathematical derivation of recall, precision, and $F_1$ scores.
  \item \textbf{Chapter~6 (HPC Implementation)}: Analyzes the high-performance computing implementation deployed on the CINECA Leonardo supercomputer (EuroHPC Booster partition), detailing SLURM job orchestration, MPI and GPU process allocation, parallel I/O throughput, and execution scaling profiles.
  \item \textbf{Chapter~7 (Results)}: Delivers a comprehensive empirical evaluation of the benchmarking campaign on the Northeastern Italy dataset, presenting pick detection recall curves, timing residual distributions, hypocentral relocation offsets, magnitude fidelity, and comparative associator analysis.
  \item \textbf{Chapter~8 (Conclusion)}: Synthesizes the core scientific and engineering findings, discusses the implications for modern operational seismological observatories, outlines current limitations, and charts future directions for AI-driven seismological monitoring.
  \item \textbf{Data and Resources}: Declares the persistent digital object identifiers (DOIs), software repositories, and access endpoints for all code, models, and catalogs developed in this work.
  \item \textbf{Appendix}: Documents supplementary mathematical derivations, configuration schemas, and detailed per-station performance tables.
\end{itemize}
